\begin{textsample}{POS Dim 1 – pi – Score 40.00 – pi/pi\_em\_40.txt}  \label{ex:f1_pos_086}
Área de Matemática e suas Tecnologias Parte Comum do Currículo Alinhado à BNCC e a DCN do Ensino Médio Nos últimos anos do \textbf{século} XIX, a educação passou por grandes transformações.
Estas transformações se deram por grandes movimentos, como, por exemplo, a Escola Nova e a Escola Ativa, \textbf{que} vieram dar às instituições de ensino \textbf{um} viés transformador, mudando o seu modo de ser através de \textbf{métodos} ativos, onde integrar o aluno ao ato de “ fazer ” se torna elemento essencial no processo de ensino e aprendizagem.
Muito antes disso, na Antiguidade, as pessoas se reuniam em várias situações para conversar, discutir, trocar ideias, e sem perceberem, \textbf{umas} ensinavam aquilo \textbf{que} sabiam às outras, de forma prática, experimental e investigativa.
E embora tais \textbf{métodos} e movimentos tenham surgido para traçar \textbf{um} novo painel ao ensino propedêutico, o \textbf{que} se \textbf{vê} na maioria das escolas é o distanciamento cada vez maior entre o \textbf{que} se deve ensinar e o \textbf{que} se deve aprender, tornando a educação escolar enfadonha e desgastada, segundo a visão dos estudantes.
A Matemática, \textbf{uma} das mais questionadas pela metodologia pouco atrativa, vem sendo ensinada sem se relacionar com a realidade e expectativa dos alunos, não \textbf{instigando} nenhuma curiosidade naquilo \textbf{que} o professor \textbf{diz} estar ensinando, isto é, o ensino da \textbf{disciplina} afasta os estudantes, ao invés de \textbf{instigar} a paixão pelos números.
Em suma, não estamos preparando nossos jovens e adultos para a revolução digital, pois o futuro será daqueles \textbf{que} saibam Matemática suficiente para se integrar nesse processo.
Infelizmente, a manipulação de poder ou a \textbf{abordagem} didática, muitas vezes exercida pelo professor durante a aula, conduz claramente à não \textbf{percepção} do \textbf{que} os silêncios na sala de aula indicam, a ausência do diálogo no contexto escolar, e isso é \textbf{um} dos motivos do fracasso do ensino de Matemática.
Segundo Baraldi, “ para os alunos, a Matemática consiste num manipular de fórmulas \textbf{que}, após certo treino, torna -se fácil em situações próprias da Matemática ” ( Baraldi, 1999, p. 88 ).
No entanto, a Matemática não é \textbf{um} cardápio de fórmulas em \textbf{que} se tenha \textbf{que} decorá -las para resolver \textbf{um} problema, mas, sobretudo, \textbf{uma} forma de pensar e enxergar o mundo \textbf{que} nos auxilia no dia a dia. \textbf{Educar} em Matemática requer objetivos estabelecidos em função de objetos do conhecimento, planejamento para a ação educativa e ferramentas para potencializar o ensino e, não menos importante, a avaliação daquilo \textbf{que} foi realizado.
É \textbf{preciso} eliminar a concepção dos \textbf{que} pensam a sala de aula como lugar de encontro de alunos ignorantes, \textbf{que} têm ali \textbf{um} professor totalmente sábio repassando informações puramente técnicas, sem vida própria e distante da realidade.
É notório \textbf{que} compreendamos a sala de aula como \textbf{um} espaço de múltiplas interações de pessoas com conhecimentos de senso comum e \textbf{que} ali estão almejando a aquisição \textbf{sistematizada} de conhecimentos.
Para Chamay ( 2001, p. 36 ), “ \textbf{um} dos objetivos essenciais do ensino da Matemática é precisamente \textbf{que}, o \textbf{que} se ensine, esteja carregado de significado, tenha sentido para o aluno ”.
O sentido de \textbf{um} conhecimento matemático se define não só pela coleção de situações em \textbf{que} este conhecimento é realizado como \textbf{teoria} matemática ou como situações em \textbf{que} o \textbf{sujeito} o encontrou como meio de solução, mas também pelo conjunto de concepções \textbf{que} rejeita, de \textbf{erros} \textbf{que} evita, de economias \textbf{que} procura, de formulações \textbf{que} retorna ( Bousseau, 1983 ).
O objeto do conhecimento matemático, quando compreendido, tem por significado inicial ter a consciência de \textbf{que} ele tem, em algum lugar, \textbf{um} tipo de existência, e \textbf{que} tal existência é algo dependente de aspectos históricos e sociais.
Assim, a Matemática deve ser inicialmente observada como ela se apresenta, levando à compreensão de \textbf{que} a investigação e a descoberta, ao longo dos anos, se constituíram como essenciais a ela, bem como seus núcleos inspiradores : \textbf{um} mundo externo com seus problemas e sua própria estruturação interna.
Sobre isso, Cortella ( 1998, p. 102 ) afirma \textbf{que} : > Quando \textbf{um} educador(a) nega ( com ou sem intenção ) aos alunos a compreensão das condições culturais, históricas e sociais de produção do conhecimento, termina por reforçar a mitificação e a sensação de perplexidade, impotência e incapacidade cognitiva.
O ensino proposto pela escola do \textbf{século} XXI desafia formar pessoas solidárias, criativas e autônomas, capazes de lidar em meio às incertezas, buscando alcançar \textbf{uma} sociedade bem mais \textbf{justa}, digna e solidária.
O ensino de Matemática é bem maior do \textbf{que} apenas provocar o \textbf{raciocínio} lógico-dedutivo dos alunos, \textbf{que} por muitas vezes parece ser o objetivo principal em planejamentos de aulas, pois através do conhecimento da história da Matemática é possível a compreensão do presente, o entendimento do passado e a projeção de futuro.
O Lúdico também é parte importante nesse processo, conectando o ensino com experiências adequadas a cada faixa etária.
Além disso, o uso de tecnologias se incorpora ao processo cognitivo como parte indissociável e indispensável para a formação na sociedade da Comunicação e da Tecnologia da atualidade.
Segundo os Parâmetros Curriculares Nacionais ( 1997 ) : > ( … ) A Matemática é componente importante na construção da cidadania, na medida em \textbf{que} a sociedade utiliza, cada vez mais, de conhecimentos científicos e recursos tecnológicos, dos quais os cidadãos devem se apropriar.
A aprendizagem em Matemática está ligada à compreensão, isto é, à apreensão do significado ; aprender o significado de \textbf{um} objeto ou acontecimento pressupõe \textbf{vê} -lo em suas relações com outros objetos e acontecimentos.
Recursos didáticos como jogos, livros, vídeos, calculadora, computadores e outros materiais têm \textbf{um} papel importante no processo de ensino aprendizagem. \textbf{Contudo}, eles precisam estar integrados a situações \textbf{que} levem ao exercício da análise e da reflexão, em última \textbf{instância}, a base da atividade matemática.
Com a BNCC, a Matemática para o Ensino Médio é apresentada com foco na construção de \textbf{uma} visão integrada, aplicada à realidade, nos diversos contextos, levando em consideração as vivências cotidianas dos \textbf{discentes}, estes \textbf{que} são inteiramente impactados pelas revoluções tecnológicas, pelo \textbf{que} \textbf{exige} o mercado de trabalho, pelos projetos de vida dos povos, entre outros.
Assim, a área da Matemática e suas Tecnologias no Ensino Médio tem a importante responsabilidade de aproveitar os conhecimentos adquiridos pelos estudantes durante o Ensino Fundamental, possibilitando o avanço no letramento matemático aprendido nos anos anteriores de escolaridade.
Segundo a nova proposta, aprender Matemática deve consistir, essencialmente, em fazer Matemática.
E fazer Matemática na escola significa tratar com questões \textbf{que} \textbf{exigem} observação, análise, estabelecimento de conexões, conjecturas, \textbf{percepção} e expressão de regularidades, busca de explicações, criação de soluções, invenção de estratégias próprias \textbf{que} envolvam noções, conceitos e \textbf{métodos} matemáticos e, por fim, a comunicação da produção realizada.
Argumentação e \textbf{prova} são dois aspectos da comunicação em Matemática.
Na sociedade \textbf{moderna} e complexa \textbf{que} \textbf{vivemos} \textbf{hoje}, ensinar Matemática requer o desenvolvimento de habilidades básicas, através de \textbf{uma} aprendizagem significativa e \textbf{que} evidencie sua \textbf{aplicabilidade}, \textbf{que} dê as possibilidades reais de conquistas.
A Educação Matemática deve visar a construção de \textbf{um} saber para capacitar nossos alunos a pensar e a refletir sobre a realidade, assim como a agir e transformá -la.
Santaló ( 2001, pg. 11 ) afirma ainda \textbf{que} : > A missão dos educadores é, portanto, preparar as novas gerações para o mundo em \textbf{que} terão \textbf{que} \textbf{viver}.
Isto quer \textbf{dizer} proporcionar -lhes o ensino necessário para \textbf{que} adquiram as destrezas e habilidades \textbf{que} vão \textbf{necessitar} para seu desempenho, com comodidade e eficiência, no seio da sociedade \textbf{que} enfrentarão ao concluir sua escolaridade.
Nesse aspecto, o ensino de Matemática prestará sua contribuição à medida \textbf{que} forem exploradas metodologias \textbf{que} priorizem a criação de estratégias, a comprovação, a justificativa, a argumentação, o \textbf{espírito} crítico, e favoreçam a criatividade, o trabalho coletivo, a iniciativa pessoal e a autonomia advinda do desenvolvimento da confiança na própria capacidade de conhecer e enfrentar desafios ( PCN’s – Matemática, 1997 ). \textbf{Enquanto} isso, a melhoria da qualidade da educação escolar no Brasil se firma, na atualidade, como \textbf{um} dos grandes desafios a serem superados pelas gestões públicas e \textbf{pesquisadores} da educação, dada a condição pouco satisfatória da qualidade da aprendizagem apresentada pela maioria dos estudantes ao final da Educação Básica, sejam estes oriundos de escolas públicas ou privadas.
A \textbf{prova} do PISA – Programa Internacional de Avaliação de Estudantes, coordenada pela OCDE – Organização para a Cooperação e o Desenvolvimento Econômico, da qual o Brasil participa como convidado desde 2000, demonstra a real situação do país em \textbf{termos} de aprendizagem, já \textbf{que} tem o objetivo de avaliar de forma amostral se os estudantes de escolas públicas e particulares, aos 15 anos de idade, adquiriram conhecimentos e habilidades essenciais para \textbf{uma} participação plena em sociedade \textbf{moderna}.
Ao avaliar a proficiência dos estudantes de 15 anos de idade em Leitura, Matemática e Ciências de escolas públicas e privadas, \textbf{que} estejam matriculados a partir do 8º ano do Ensino Fundamental, o objetivo principal do PISA é fornecer indicadores \textbf{que} contribuam para a discussão da qualidade da educação ofertada nos países participantes, de modo a subsidiar políticas de melhoria da educação básica.
Por outro \textbf{lado}, os resultados produzidos pelo programa podem indicar até \textbf{que} ponto as escolas de cada país participante estão preparando seus jovens para exercer e/ou responder às demandas sociais, já \textbf{que} o PISA avalia até \textbf{que} ponto os estudantes, próximos ao final da educação obrigatória, adquiriram conhecimentos e habilidades essenciais para plena participação na vida social e econômica.
Conforme Pessoa e Silva ( 2019 ), o resultado brasileiro nesta avaliação comparada, isto é, promovida entre os países e convidados da OCDE, indica \textbf{que} os estudantes brasileiros avaliados em 2012, cujo foco da avaliação era Matemática, possuem baixos níveis de aprendizado matemático, considerando \textbf{que} cerca de 70 \% encontram -se nos níveis menores da escala, conforme a seguir : | Nível/Edição | Abaixo de 1 | 1 | 2 | 3 | 4 | 5 | 6 | |---|---:|---:|---:|---:|---:|---:|---:| | 2012 | 35 \% | 32 \% | 21 \% | 8 \% | 4 \% | 1 \% | 0,5 \% | Tabela 1 – Resultado Brasil 2012 – Matemática Fonte : INEP ( 2013 ).
Para as \textbf{autoras}, esta constatação é validada ao observarmos a diferença entre o Brasil e países \textbf{que} obtêm os melhores resultados no PISA, verificada nas etapas posteriores da vida escolar.
Segundo a OCDE, a média de engenheiros formados entre os concluintes de cursos superiores nos países membros da organização é de 14 \% ; na Coreia do Sul, esse índice é de 25 \%.
No Brasil, segundo dados do MEC/INEP, obtidos através do Censo da Educação Superior ( 2016 ), apenas cerca de 6 \% dos concluintes estavam nas áreas de Engenharia.
A distância observada na condução das políticas educacionais é ainda mais esmagadora se considerarmos o desenvolvimento tecnológico e econômico observado entre os dois países.
Nesse ínterim, é, portanto, válido considerar através desses percentuais a vocação e o incentivo \textbf{que} os dois países fornecem para a inovação tecnológica.
Segundo o boletim Radar : tecnologia, produção e comércio exterior do Instituto de Pesquisa Econômica e Aplicada – IPEA ( 2014 ) : > [ … ] a economia nacional mostra -se pouco intensiva em trabalho qualificado e de cunho técnico-científico, insumos tidos como necessários à inovação e à competição em mercados globais.
De outro, o sistema educacional brasileiro parece \textbf{operar} em \textbf{um} equilíbrio de baixa qualidade, em todos os níveis de ensino, o \textbf{que} afeta a capacidade produtiva da força de trabalho, em geral, e a das engenharias, em particular ( RADAR, 2014, p. 19 ).
A formação em nível superior nas carreiras \textbf{que} \textbf{exigem} aprofundamento científico fica comprometida pela pouca eficiência dos sistemas de ensino ofertados nacionalmente e, nesse âmbito, “ o Brasil segue tendo diante de si grandes desafios para melhorar o aprendizado, em todos os níveis de escolaridade, com potenciais repercussões na qualidade da força de trabalho ” ( RADAR, p. 34 ).
Em âmbito estadual, o Piauí não foge à regra brasileira : desde 2006 recebendo nota padronizada no teste, o Estado obteve resultados mais baixos \textbf{que} a média nacional, embora com pequenos avanços.
Conforme dados do PISA/INEP ( 2016 ), o Piauí concentra a maioria dos alunos ( 79 \% ) nos níveis abaixo de 1 e 1 em Matemática.
Tais números \textbf{revelam} \textbf{que} o padrão de desempenho dos estudantes piauienses demonstra carência de aprendizagem em relação ao \textbf{que} é previsto para a etapa de escolaridade avaliada, tendo em vista \textbf{que} a OCDE institui o nível 2 como o “ aceitável ” para os estudantes com 15 anos de idade.
Dessa forma, o resultado do Piauí na avaliação corrobora a fragilidade do sistema educacional brasileiro diante das desigualdades econômicas e sociais \textbf{historicamente} verificadas e debatidas em âmbito nacional, se incorporando às diferentes concepções educacionais observadas nas regiões brasileiras.
Tradicionalmente, os estados da região Norte e Nordeste detêm, em média, os piores indicadores educacionais e sociais ; no PISA, essa realidade não é diferente.
Tais desigualdades educacionais estão muito presentes no cenário brasileiro.
Na tentativa de amenizar essas irregularidades, a Base Nacional Comum Curricular – BNCC se consolida como \textbf{uma} oportunidade para a transformação dos \textbf{sujeitos}, a partir da aquisição de novas competências.
A reformulação curricular proposta com a BNCC tem o objetivo de \textbf{aproximar} mais a escola das expectativas de aprendizagem dos estudantes do \textbf{século} XXI, \textbf{que} são altamente tecnológicos e \textbf{que} buscam respostas para suas demandas socioemocionais, evidenciando o protagonismo juvenil/estudantil, bem como sua preparação para o mundo do trabalho.

% matched lemmas: abordagem, aplicabilidade, aproximar, autor, contudo, discente, disciplina, dizer, educar, enquanto, erro, espírito, exigir, historicamente, hoje, instigar, instância, justo, lado, moderno, método, necessitar, operar, percepção, pesquisador, preciso, prova, que, raciocínio, revelar, sistematizar, sujeito, século, teoria, termos, um, ver, vivar, viver
\end{textsample}
